% ----------------------------------------------------------
% Introdução
% ----------------------------------------------------------
\chapter{Introdução}\label{cap:introducao}
%\addcontentsline{toc}{chapter}{Introdução}
% ----------------------------------------------------------

A internet cresce dia após dia e cada vez mais diferentes tipos de dispositivos são conectados a essa imensa rede. Há uma estimativa de 26 bilhões de ''coisas'' conectadas a internet até 2020, comparado a 6 bilhões na década de 2000. Isso foi possível graças ao custo do acesso a internet e banda larga terem diminuído 40 vezes nos últimos 10 anos. Esse \textit{boom} de crescimento também é influenciado pela IoT (\textit{Internet of Things} - Internet das Coisas). \cite{goldmansachs-iot}. 

Segundo \citeonline{ashton-iot} o termo IoT surgiu como título de sua apresentação à Procter \& Gamble (P\&G) em 1999. Esse termo advém do uso de internet em diferentes tipos de dispositivos. 

\begin{citacao}
''(...) a Internet das Coisas é um conceito no qual dispositivos de nosso dia a dia são equipados com sensores capazes de captar aspectos do mundo real, como por exemplo, temperatura, umidade, presença, etc, e envia-los a centrais que recebem estas informações e as utilizam de forma inteligente.'' \cite{nascimento-iot}.
\end{citacao}

Como exemplo podemos citar: geladeira ligada a internet informando a falta de condimentos, caixa de remédio conectada a internet prevendo o término de remédio para avisar ao consumidor, entre diversos outros. Um bom exemplo é interligar uma casa por meio de sensores, como termostatos, sistemas de segurança, iluminação, sistemas de entretenimento com uma inteligência por trás para diversas aplicações. \cite{goldmansachs-iot}

Diversas áreas tiveram o seu crescimento alavancado por conta da IoT. O mais marcante é o ramo do DIY (\textit{Do It Yourself}, ou faça você mesmo), em que pessoas criam ou adaptam coisas para suas necessidades. Isso se dá por conta de ter aparecido no mercado dispositivos e componentes que auxiliam a criação e adaptação de eletrônicos e outros. Um ótimo exemplo é o Arduino, um dispositivo que ajuda a criar os projetos de eletrônica e que é composto de duas partes: o hardware e o software. Com eles é possível construir praticamente de tudo, desde um LED piscante a um robô que envia um \textit{tweet} quando sua planta está sem água. \cite{ben-arduino}. \cite{sorrel-arduino}.

Além do Arduino existem diversos outros \textit{devices} com funcionalidades similares ou até complementares. Podemos citar o \textit{Raspberry Pi}, um computador do tamanho de um cartão de crédito e de baixo custo. \cite{raspberrypi-rpi}. Por ser um computador é possível executar um sistema operacional (como Linux, \textit{RISC OS}, \textit{Windows 10 for IoT}). Dependendo do modelo possui portas USB (1 a 4) para conexão com periféricos (como mouse, teclado, adaptador WiFi, Bluetooth), e também porta Ethernet para conexão a internet cabeada (exceto modelo A e A+). Também possui de 26 (modelo A e B) a 40 pinos (modelo A+, B+ e 2) para conexões gerais de entrada e saída digitais (GPIO - \textit{General Input and Output}). 

Através desses pinos pode-se conectar uma diversidade de componentes eletrônicos como sensores, atuadores, outros dispositivos para comunicação. Dessa forma, suas funcionalidades são expandidas de uma forma absurda, ficando a cargo de cada pessoa montar uma nova aplicação. Uma boa aplicação é a conexão de um adaptador bluetooth pela USB para comunicação com smartphones, tablets, PCs e outros dispositivos tipo sensores sem fio.

Um exemplo desses sensores externos são os \textit{beacons}, pequenos modelos sem fio que se comunicam por meio do Bluetooth 4.0 ou BLE (\textit{Bluetooth Low Energy} - Bluetooth de Baixa Energia). Como o próprio nome diz, esse padrão de comunicação utiliza pouca energia. Desta forma, um \textit{beacon} pode funcionar por anos. ''Na prática, ela permite localizar objetos (ou pessoas que carregam esses objetos) com muito mais precisão dentro de ambientes fechados.'' \cite{teixeira-beacon}.

Os \textit{beacons} foram pouco explorados até o momento, seu uso está sendo mais notado na área de grandes lojas do varejo.

\begin{citacao}
''A Apple (...) já está utilizando a tecnologia em 254 lojas nos EUA. As funcionalidades já estão embutidas na versão oficial do aplicativo da Apple Store para iOS [, sistema operacional de seus smartphones e tablets]. (...) Quando o usuário se aproxima de uma loja física, o aplicativo oferece toda uma camada extra de informações e serviços que são específicos para aquela unidade – como por exemplo ofertas locais, tamanho da fila para ser atendido no Genius Bar, eventos e treinamentos que estão agendados ali na loja etc.'' \cite{teixeira-beacon}.
\end{citacao}

A Macy's (grande rede norte americana de loja de departamentos) realizou testes em algumas de suas lojas para enviar alertas a pessoas que entrarem em suas lojas, com promoções e melhores indicações, utilizando o padrão \textit{iBeacon} da Apple. Os testes foram limitados a usuários de iPhone, e somente quando entraram na loja. Em um futuro teste espera-se que seja possível separar por departamentos, para que quando um usuário percorra a loja apareça as notícias relativas ao local da loja que ela está. \cite{kastrenakes-macys-beacon}


% ---
% Objetivos
% ---
\section{Objetivos}\label{cap:objetivos-justificativa}
% ---

% ---
% Objetivos Gerais
% ---
\subsection{Objetivos Gerais}\label{sec:objetivos-gerais}
% ---

O objetivo geral desse projeto é projetar e desenvolver um dispositivo rastreador de \textit{iBeacons} utilizando o \textit{Raspberry Pi}.

% ---

% ---
\subsection{Objetivos Específicos}\label{sec:objetivos-especificos}
% ---

\begin{alineas}
	\item Estudar o funcionamento de \textit{beacons}, comunicação via \textit{bluetooth low energy} e \textit{Raspberry Pi}, assim como suas aplicações;
	\item Identificar os requerimentos básicos de funcionamento de \textit{beacons} e \textit{Raspberry Pi};
	\item Definir os elementos para desenvolvimento do protótipo;
	\item Planejar a estrutura do sistema;
	\item Implementar o protótipo de acordo com a estrutura e elementos planejados.
\end{alineas}

% ---


% ---
\section{Logotipo e Nome para o Dispositivo}\label{sec:logotipo-prototipo}
% ---

Por ser um produto, decidiu-se que o dispositivo deveria ter um nome e logotipo. Durante algumas reuniões de \textit{brainstorm}, o nome \textit{Peacon} surgiu e foi o escolhido. Esse nome deriva das palavras \textit{Pi}  + \textit{beacon}. Nos próximos capítulos o dispositivo rastreador será referido conforme seu nome.

Juntando todos os atributos do Peacon, criou-se um logotipo\footnote{O logotipo foi criado por Henrique Lopes Silva Santos, aluno do curso de Design da Faculdade de Arquitetura, Artes e Comunicação da UNESP de Bauru.} que identificasse muito bem todas as suas funcionalidades (\autoref{fig:logo-peacon}). Adesivos foram criados para colar no protótipo e identificar o produto.

\begin{figure}[htb]
	\caption{\label{fig:logo-peacon}Logotipo criado para o protótipo}
	\begin{center}
		\includegraphics[width=0.5\textwidth]{img/logo-peacon.png}
	\end{center}
	\legend{Fonte: elaborado pelo autor}
\end{figure}

% ----------------------------------------------------------